% ===================================================================
%  TAULA N.º 5 — La Casa i la seva construcció.
%  Traduction 2026 d'apres texto/io, controlee sur texto/fr et
%  texto/es. Consigne : voir texto/ca/10-taula-01.tex.
%
%  Le tableau 5 est un DIALOGUE : on garde \textsc{} pour le nom du
%  parleur. Les renvois du plan sont des LETTRES, (a) a (m), et l'on
%  ecrit ce que l'ido ecrit, « (1) » compris la ou l'imprimeur a
%  compose un 1 pour un l.
%
%  LES DEUX NOMS PROPRES SE TRADUISENT, ET L'IDO L'A DEJA FAIT :
%  Leblanc est « Blanko », Roux est « Rufo ». Le catalan a « Blanc »
%  et « Roig », qui sont l'un et l'autre de vrais noms de famille.
% ===================================================================

%%K t05-apar-1 apar
\VUsaut{6.0mm}
\VUtitre{15.0pt}{60}{TAULA N.º 5}
\VUsaut{1.4mm}
\VUtitre{11.4pt}{0}{La Casa i la seva construcció.}
\VUsaut{2.2mm}

%%K t05-tit-1 sub
\VUpk{10.2pt}{40}{Una bona trobada.}
\VUsaut{3.0mm}

%%K t05-01-1 p
\textsc{Joan}. --- Oncle, m'agradaria molt saber com es construeix una
\VUgras{casa}.

%%K t05-01-2 p
\textsc{L'Oncle} \textsuperscript{(80)}. --- És molt fàcil, fill meu; vine amb mi
i et mostraré la que el senyor Bernat \textsuperscript{(79)} està fent construir i
en la qual viuré.

%%K t05-01-3 p
\textsc{Joan}. --- I qui ha dibuixat el \VUgras{plànol} \textsuperscript{(76)}
d'aquesta casa?

%%K t05-01-4 p
\textsc{L'Oncle}. --- L'\VUgras{arquitecte} \textsuperscript{(75)}; va presentar
un dibuix molt clar, que va ser aprovat, i el \VUgras{contractista}
\textsuperscript{(77)} va començar les obres.

%%K t05-01-5 p
\textsc{Joan}. --- I qui és el contractista?

%%K t05-01-6 p
\textsc{L'Oncle}. --- És el senyor Blanc, el pare del teu amic Jaume. Dirigeix
els obrers amb el senyor Roig, l'arquitecte.

%%K t05-01-7 p
Mira! Aquí arriba justament el senyor Blanc amb el seu \VUgras{regle}
\textsuperscript{(78)}, i el senyor Roig amb el seu plànol.

%%K t05-01-8 p
Bon dia, senyors. Com estan?

%%K t05-01-9 p
\textsc{El Sr. Roig}. --- Molt bé, gràcies. Bon dia, Joan; com ha crescut aquest
nen!

%%K t05-01-10 p
\textsc{L'Oncle}. --- Ja ho crec, és tot un homenet. És molt seriós; se li ha
d'explicar tot el que veu. Avui voldria saber com es construeix una casa.

%%K t05-tit-2 sub
\VUpk{10.2pt}{40}{Els obrers.}
\VUsaut{3.0mm}

%%K t05-01-11 p
\textsc{El Sr. Blanc}. --- Doncs vine amb mi, amiguet, i t'ho explicaré de
seguida, perquè m'agraden molt els nens que volen aprendre.

%%K t05-01-12 p
Veus aquell obrer que té una \VUgras{pala} \textsuperscript{(82)} a la mà: és un
\VUgras{peó} \textsuperscript{(81)}. Està obrint una \VUgras{rasa}
\textsuperscript{(46)} per col·locar-hi la canonada de l'aigua. També va ser ell
qui va cavar el sòl al lloc on és la casa, perquè el paleta aixequés les quatre
\VUgras{parets}. La part d'aquelles parets que queda sota terra s'anomena els
\VUgras{fonaments} \textsuperscript{(2)}. L'espai comprès entre els fonaments
forma el \VUgras{celler} \textsuperscript{(3)}. Aquell celler té una finestreta
anomenada \VUgras{claraboia} \textsuperscript{(5)}, una \VUgras{porta}
\textsuperscript{(4)} i una escala.

%%K t05-01-13 p
El \VUgras{paleta} \textsuperscript{(108)} va continuar aixecant les parets deixant
buits per a les \VUgras{portes} \textsuperscript{(8)} i les \VUgras{finestres}
\textsuperscript{(17)}.

%%K t05-01-14 p
\textsc{Joan}. --- Però no el veig, aquell paleta!

%%K t05-01-15 p
\textsc{El Sr. Blanc}. --- Ja ha acabat de treballar a la casa. És vora la
\VUgras{quadra} \textsuperscript{(54)}, damunt la seva \VUgras{bastida}
\textsuperscript{(110)}, aixecant la paret del \VUgras{safareig}
\textsuperscript{(56)}.

%%K t05-01-16 p
Té una plana, una gaveta i una \VUgras{plomada} \textsuperscript{(109)}. Però no
te'n recordaràs, d'aquests noms.

%%K t05-01-17 p
\textsc{Joan}. --- I tant que me'n recordaré, perquè demanaré al meu oncle una
plana petita i una gaveta, i jo mateix faré una plomada amb un cordill.

%%K t05-tit-3 sub
\VUsaut{3.0mm}
\VUpk{10.2pt}{40}{Els materials.}
\VUsaut{3.0mm}

%%K t05-01-18 p
\textsc{L'Oncle}. --- Ah, molt bé! Però digues, Joan, què cal per construir una
casa?

%%K t05-01-19 p
\textsc{Joan}. --- Calen \VUgras{carreus} \textsuperscript{(59)} i \VUgras{morter}
\textsuperscript{(65)}.

%%K t05-01-20 p
\textsc{L'Oncle}. --- Exacte. Mira aquell home del barret gran, al seu
\VUgras{carro} \textsuperscript{(136)}. És un \VUgras{carreter}
\textsuperscript{(135)}. Cada dia porta aquí \VUgras{blocs de pedra}
\textsuperscript{(60)}. Descarrega el seu carro al pati, i els \VUgras{picapedrers}
\textsuperscript{(83)} llauren els blocs i els serren amb la seva \VUgras{serra de
pedra} \textsuperscript{(85)}. Un \VUgras{manobre} \textsuperscript{(146)} els
porta al paleta, que els col·loca.

%%K t05-01-21 p
Mentrestant, el \VUgras{pastador} \textsuperscript{(87)} prepara el morter.
Necessita \VUgras{sorra} \textsuperscript{(62)}, \VUgras{calç}
\textsuperscript{(63)} i \VUgras{aigua} \textsuperscript{(64)}. La calç és al seu
costat en un \VUgras{sac} \textsuperscript{(71)}; un \VUgras{ajudant de paleta}
\textsuperscript{(111)} li porta la sorra en un \VUgras{carretó}
\textsuperscript{(112)}, i l'\VUgras{aprenent} \textsuperscript{(113)} és a la
bomba \textsuperscript{(58)}. Omple un \VUgras{cubell}
\textsuperscript{(114)}. El morter estarà a punt de seguida. El \VUgras{portador
de morter} \textsuperscript{(107)} l'agafa i el puja per l'escala fins a la
bastida, on un paleta acaba aquell costat de la casa.

%%K t05-01-22 p
I ara, com s'anomena l'obrer que fa la \VUgras{teulada} \textsuperscript{(31)}?

%%K t05-01-23 p
\textsc{Joan}. --- És el \VUgras{teulader} \textsuperscript{(118)}.

%%K t05-tit-4 sub
\VUsaut{3.0mm}
\VUpk{10.2pt}{40}{La teulada de la casa.}
\VUsaut{3.0mm}

%%K t05-01-24 p
\textsc{El Sr. Blanc}. --- Sí, però primer ve el \VUgras{fuster de gruixa}
\textsuperscript{(115)}.

%%K t05-01-25 p
El fuster de gruixa fa l'\VUgras{encavallada} \textsuperscript{(35)}. Uneix les
\VUgras{bigues} \textsuperscript{(37)} amb \VUgras{clavilles}
\textsuperscript{(117)}. Aquells obrers de dalt, amb els seus amples calçons de
vellut, són fusters de gruixa. Un sosté una biga petita i l'altre talla una
clavilla amb la seva \VUgras{destral} \textsuperscript{(116)}. El qui és vora la
\VUgras{xemeneia} \textsuperscript{(41)} és el teulader. Clava llistons damunt les
\VUgras{bigues} (*), i \VUgras{pissarres} \textsuperscript{(66)} damunt els
llistons. De vegades hi posa teules.

%%K t05-noto-1 noto
\VUnotes{143.2mm}{%
(*) \VUgras{Balk-o} = alemany \textit{Balken}, anglès \textit{joist}, francès
\textit{solive}, italià \textit{travicello}, rus \textit{balka}, espanyol i
portuguès \textit{viga}. Arrel tècnica encara no adoptada, però que es proposa
aquí per traduir el sentit del mot francès \textit{solive}, que no és una
\textit{jàssera} (francès \textit{poutre}) ni una biga petita. És un fust
escairat que sosté un sòl i un sostre.%
}

%%K t05-01-26 p
La teulada de la quadra és de \VUgras{teules} \textsuperscript{(55)}, la de la
casa és de \VUgras{pissarra} \textsuperscript{(32)} i la de la galeria és de
\VUgras{zinc} \textsuperscript{(24)}.

%%K t05-01-27 p
\textsc{Joan}. --- El teulader fa també les teulades de zinc?

%%K t05-01-28 p
\textsc{El Sr. Blanc}. --- No, aquest és el \VUgras{zinquer}
\textsuperscript{(121)} o el \VUgras{llauner}, el qui està col·locant una
\VUgras{golfa} \textsuperscript{(38)} a la teulada de la casa. També posa els
\VUgras{canals} \textsuperscript{(43)} i els \VUgras{baixants}
\textsuperscript{(42)}. Solda el zinc i la llauna. Va ser ell qui va fixar el
\VUgras{parallamps} \textsuperscript{(40)} i el \VUgras{penell}
\textsuperscript{(39)} al \VUgras{frontis} \textsuperscript{(36)}.

%%K t05-01-29 p
\textsc{Joan}. --- Aleshores la casa està acabada?

%%K t05-tit-5 sub
\VUsaut{3.0mm}
\VUpk{10.2pt}{40}{Les eines.}
\VUsaut{3.0mm}

%%K t05-01-30 p
\textsc{El Sr. Blanc}. --- No del tot. El \VUgras{guixaire}
\textsuperscript{(99)} aixeca amb \VUgras{maons} \textsuperscript{(61)} els envans
que la divideixen en les diferents habitacions. Revesteix de \VUgras{guix}
\textsuperscript{(70)} els \VUgras{sostres} \textsuperscript{(26)} i les parets.
Ara està fent el sostre de la \VUgras{galeria} \textsuperscript{(33)}. Té, com el
paleta, una \VUgras{plana} \textsuperscript{(100)} i una \VUgras{gaveta}
\textsuperscript{(101)}.

%%K t05-01-31 p
Després l'ebenista fa les portes, les finestres, els \VUgras{parquets}
\textsuperscript{(16)} i els \VUgras{porticons} \textsuperscript{(19)}.

%%K t05-01-32 p
Ara treballa al pati al seu \VUgras{banc de fuster} \textsuperscript{(90)}. Té una
\VUgras{serra} \textsuperscript{(94)} per serrar la \VUgras{fusta}
\textsuperscript{(69)}, un \VUgras{ribot} \textsuperscript{(91)}, un
\VUgras{serjant} \textsuperscript{(92)}, un \VUgras{enformador}
\textsuperscript{(94 \textit{bis})}, un \VUgras{compàs} \textsuperscript{(95)} i
una \VUgras{maça de fusta} \textsuperscript{(95 \textit{bis})}.

%%K t05-01-33 p
\textsc{Joan}. --- Ah, però encara és més divertit ser ebenista que paleta!
Oncle, em compraràs un banc, una serra i un ribot? Faré una caseta per a Medor.

%%K t05-01-34 p
\textsc{L'Oncle}. --- Sí, fill meu.

%%K t05-01-35 p
\textsc{Joan}. --- Que content que estic! I aquell que és vora la porta
d'entrada, qui és?

%%K t05-tit-6 sub
\VUsaut{3.0mm}
\VUpk{10.2pt}{40}{La pintura.}
\VUsaut{3.0mm}

%%K t05-01-36 p
\textsc{El Sr. Blanc}. --- És un \VUgras{pintor} \textsuperscript{(102)}. És a la
seva escala amb el seu pot de \VUgras{pintura} \textsuperscript{(103)} i la seva
\VUgras{brotxa} \textsuperscript{(104)}. Pinta la fusta de la porta d'entrada
perquè es conservi més temps.

%%K t05-01-37 p
També és ell qui enganxa el paper a les habitacions.

%%K t05-01-38 p
El \VUgras{tapisser} \textsuperscript{(107)} que és en una \VUgras{escala}
\textsuperscript{(106)}, al \VUgras{balcó} \textsuperscript{(28)}, col·loca un
\VUgras{tendal} \textsuperscript{(29)}.

%%K t05-01-39 p
L'obrer que és vora la finestra gran és el \VUgras{vidrier}
\textsuperscript{(96)}. Fixa els \VUgras{vidres} \textsuperscript{(18)} amb
\VUgras{massilla} \textsuperscript{(73)}. Aquell altre obrer, agenollat vora la
porta del celler, és un \VUgras{serraller} \textsuperscript{(97)}. Està posant un
\VUgras{pany} \textsuperscript{(6)}. La seva \VUgras{llima} \textsuperscript{(98)}
és al seu costat.

%%K t05-01-40 p
\textsc{Joan}. --- És també serraller aquell que colpeja amb el seu
\VUgras{martell} \textsuperscript{(84)} damunt un ferro?

%%K t05-01-41 p
\textsc{El Sr. Blanc}. --- Més aviat és un \VUgras{ferrer}
\textsuperscript{(125)}. Forja \VUgras{ferros} \textsuperscript{(67)} per a
l'\VUgras{electricista} \textsuperscript{(124)}, que és a la teulada de la
\VUgras{cotxera} \textsuperscript{(51)}. Té una \VUgras{fornal}
\textsuperscript{(126)}, una \VUgras{enclusa} \textsuperscript{(127)} i unes
\VUgras{tenalles} \textsuperscript{(128)}.

%%K t05-01-42 p
L'electricista col·loca el \VUgras{fil de coure} \textsuperscript{(44)} del
telèfon.

%%K t05-tit-7 sub
\VUpk{10.2pt}{40}{La jardineria.}
\VUsaut{3.0mm}

%%K t05-01-43 p
\textsc{L'Oncle}. --- Al fons del \VUgras{pati} \textsuperscript{(45)}, vora la
\VUgras{reixa} \textsuperscript{(47)} i el \VUgras{portal}
\textsuperscript{(48)}, veus el \VUgras{jardiner} \textsuperscript{(129)}. Està
preparant el petit \VUgras{jardí} \textsuperscript{(49)}. Ha cavat la terra amb
la seva \VUgras{aixada} \textsuperscript{(131)}, sembra llavors i rastella amb el
\VUgras{rasclet} \textsuperscript{(130)}. Després, amb la seva \VUgras{regadora}
\textsuperscript{(132)}, regarà els \VUgras{planters} \textsuperscript{(150)} i
les plantes creixeran.

%%K t05-01-44 p
El \VUgras{mosso de quadra} \textsuperscript{(133)}, que acaba de cuidar el
cavall i de donar-li menjar, neteja el \VUgras{cotxe} \textsuperscript{(52)} amb
una esponja, una \VUgras{bomba de mà} \textsuperscript{(134)} i un cubell
d'aigua. Després el ficarà a la cotxera i tancarà la porta amb el gruixut
\VUgras{forrellat} \textsuperscript{(53)}.

%%K t05-01-45 p
Ja estàs content, suposo; has tingut prou explicacions.

%%K t05-01-46 p
\textsc{Joan}. --- Sí, oncle, però m'agradaria veure l'interior de la casa.

%%K t05-01-47 p
\textsc{L'Oncle}. --- Això serà demà. El senyor Roig t'explicarà el plànol i et
dirà el nom de cada habitació.

%%K t05-tit-8 sub
\VUsaut{3.0mm}
\VUpk{10.2pt}{40}{La distribució interior de la casa.}
\VUsaut{2.4mm}

%%K t05-apar-2 apar
{\centering\textit{(Vegeu el plànol.)}\par}
\VUsaut{2.4mm}

%%K t05-01-48 p
\textsc{L'Arquitecte}. --- Vine, Joan, et faré visitar les diferents parts de la
casa.

%%K t05-01-49 p
Entrem a la \VUgras{planta baixa} \textsuperscript{(7)}. Heus aquí el botó del
\VUgras{timbre} \textsuperscript{(11)}. Pugem els tres \VUgras{graons}
\textsuperscript{(14)} de l'\VUgras{escalinata} \textsuperscript{(9)}, passem el
\VUgras{llindar} \textsuperscript{(10)} de la \VUgras{porta d'entrada}
\textsuperscript{(8)}, i heus-nos aquí al peu de l'\VUgras{escala}
\textsuperscript{(13)}.

%%K t05-01-50 p
\textsc{Joan}. --- Això és el passadís.

%%K t05-01-51 p
\textsc{L'Arquitecte}. --- No, això és el \VUgras{vestíbul}
\textsuperscript{(12)}. El \VUgras{passadís} \textsuperscript{(a)} és al fons. A
la dreta, heus aquí la \VUgras{sala} \textsuperscript{(b)} i el
\VUgras{menjador} \textsuperscript{(c)}; davant del menjador hi ha la
\VUgras{cuina} \textsuperscript{(d)} i el \VUgras{rebost}
\textsuperscript{(e)}; després, al davant, el \VUgras{despatx}
\textsuperscript{(f)}. La \VUgras{galeria} \textsuperscript{(23)}, amb la seva
\VUgras{porta envidrada} \textsuperscript{(25)}, és vora la sala.

%%K t05-01-52 p
Ara pujarem l'escala. Agafa't al \VUgras{passamà} \textsuperscript{(15)} i compta
els graons. Són catorze. Heus aquí el \VUgras{replà} \textsuperscript{(m)}: som al
primer \VUgras{pis} \textsuperscript{(27)}.

%%K t05-tit-9 sub
\VUsaut{3.0mm}
\VUpk{10.2pt}{40}{El primer pis.}
\VUsaut{3.0mm}

%%K t05-01-53 p
Davant de l'escala hi ha el \VUgras{comú} \textsuperscript{(20)} i el
\VUgras{tocador} \textsuperscript{(j)}. A la dreta, un bell \VUgras{dormitori}
\textsuperscript{(g)} amb dues finestres a la \VUgras{façana}
\textsuperscript{(1)} de la casa. A l'esquerra hi ha el \VUgras{bany}
\textsuperscript{(1)} i la \VUgras{cambra d'hostes} \textsuperscript{(i)} amb una
gran \VUgras{alcova} \textsuperscript{(h)}. Vora el dormitori hi ha la
\VUgras{cambra dels nens} \textsuperscript{(k)}. Dóna a una àmplia
\VUgras{terrassa} \textsuperscript{(21)}. Sota aquella terrassa hi ha un altre
comú \textsuperscript{(20)} i, contra la paret, heus aquí la \VUgras{campana}
\textsuperscript{(22)} que crida a sopar.

%%K t05-01-54 p
\textsc{Joan}. --- Pugem al \VUgras{segon pis}.

%%K t05-01-55 p
\textsc{L'Arquitecte}. --- No hi ha segon pis. Sota la teulada només hi ha
\VUgras{golfes} \textsuperscript{(33)} i les golfes altes
\textsuperscript{(34)}. Hem acabat la visita, podem baixar.

%%K t05-01-56 p
\textsc{Joan}. --- Gràcies, senyor, ha estat molt interessant. Aniré a explicar
al meu pare tot el que he vist.
\VUsaut{4.0mm}
\VUfilet{22mm}
